\clearpage
\newpage
\section*{Abstract}

The rapid advancement of artificial intelligence and the increasing
complexity of deep neural networks (DNNs) have increased the demand
for specialized hardware accelerators capable of executing these
workloads efficiently. However, designing a DNN accelerator remains
a complex task due to the large number of architectural and
implementation decisions involved. Exploring combinations of
processing element arrays, memory hierarchies, dataflows, and
hardware components under power, performance, and area (PPA)
constraints results in a large design space that is difficult to
explore manually.

This thesis presents \textbf{TALOS}, a framework for the hierarchical exploration
and generation of DNN hardware accelerators. TALOS separates the
design process into two levels. At the architecture level, candidate
accelerators are explored for a given DNN workload, considering
parameters such as the processing element array, memory hierarchy,
and algorithm-to-hardware mapping. At the implementation level,
promising architectures are further evaluated and instantiated using
characterized hardware building blocks, providing a more detailed
estimation of their PPA characteristics.

This separation allows architectural decisions to be explored
independently from specific hardware implementations, while still
providing a path from a DNN workload to a realizable accelerator
design. The framework is designed to support different optimization
objectives and hardware component libraries, allowing the exploration
flow to be adapted to different workloads and design constraints.

TALOS is evaluated using representative DNN workloads and accelerator
architectures, assessing its ability to explore the design space and
identify hardware configurations according to user-defined PPA
objectives. By combining architecture-level exploration with
hardware-aware implementation decisions, TALOS provides a modular
flow for reducing the effort required to design domain-specific DNN
accelerators.